% !TEX root = main.tex

% =============================
\documentclass[12pt,a4paper,fleqn]{report}


% --- Font (Times New Roman asli) ---
\usepackage{fontspec}
\setmainfont{Times New Roman}[Ligatures=TeX,Renderer=Basic]

\usepackage[american,indonesian]{babel}

\usepackage{csquotes}
\usepackage[style=apa,backend=biber]{biblatex}
\addbibresource{references.bib}

% --- MAP bahasa indonesian ke mapping 'american-apa' yang tersedia ---
\DeclareLanguageMapping{indonesian}{american-apa}

% --- Optional: override beberapa label bibliografi agar muncul Bahasa Indonesia ---
\DefineBibliographyStrings{indonesian}{%
  bibliography = {DAFTAR PUSTAKA},
  references = {DAFTAR PUSTAKA},
  and          = {\&},       % Mengubah "dan" menjadi "&"
  %andothers    = {dkk\adddot} % Opsional: mengubah "et al." jadi "dkk."
  %andothers = {dkk\adddot} % contoh: "et al." -> "dkk."
}

% Pilih bahasa tampilan dokumen jadi indonesian (tidak mengubah mapping di atas)
\AtBeginDocument{\selectlanguage{indonesian}}

% --- Margin ---
\usepackage{geometry}
\geometry{
  a4paper,
  left=4cm,
  right=3cm,
  top=4cm,
  bottom=3cm
}

% --- Spasi 1.5 ---
\usepackage{setspace}
\onehalfspacing

% --- Paragraf (indentasi) ---
\usepackage{indentfirst}
\setlength{\parindent}{1cm}
\setlength{\parskip}{0pt}

% --- Matematika & fisika ---
\usepackage{amsmath,amssymb,amsfonts,mathtools}
\setlength{\mathindent}{\parindent}
\usepackage{physics}
\usepackage{braket}
\usepackage{bm}
\usepackage{tensor}
\usepackage{siunitx}

% --- Gambar & grafik ---
\usepackage{graphicx}
\usepackage{caption}
\usepackage{subcaption}
\usepackage{float}
\usepackage{tikz}
\usepackage{pgfplots}
\pgfplotsset{compat=1.18}

\usepackage{booktabs}
\usepackage{array}
\usepackage{makecell}
\usepackage{colortbl}
\usepackage{tabularx}

% --- Hyperlink & referensi ---
\usepackage[hidelinks]{hyperref}
\usepackage{cleveref}

% --- Penomoran persamaan per bab ---
\numberwithin{equation}{chapter}

% --- Enumerate custom ---
\usepackage{enumitem}


% --- Format judul bab: "BAB I" dan JUDUL di tengah, nomor bab Romawi ---

% --- Pastikan titlesec dimuat ---
\usepackage{titlesec}

% --- Tetap gunakan numbering chapter ARABIC untuk penomoran internal ---
\renewcommand{\thechapter}{\arabic{chapter}}
\renewcommand{\thesection}{\thechapter.\arabic{section}}
\renewcommand{\thesubsection}{\thesection.\arabic{subsection}}

% --- Tampilakan label "BAB I" dengan ROMAWI (hanya tampilan) ---
\titleformat{\chapter}[display]
  {\normalfont\centering\bfseries\normalsize}   % ukuran 12pt (normalsize dalam dokclass 12pt)
  {\MakeUppercase{BAB \Roman{chapter}}} % CETAKAN label pakai ROMAWI tanpa mengganti \thechapter
  {0pt}
  {\vspace{1ex}\MakeUppercase}  % membuat judul bab juga kapital seluruhnya; hapus \MakeUppercase jika tidak mau
\titlespacing*{\chapter}{0pt}{0pt}{20pt}

% --- Section / subsection tetap 12pt ---
\titleformat{\section}
  {\normalfont\bfseries\normalsize}
  {\thesection}{1em}{}
\titlespacing*{\section}{0pt}{12pt}{6pt}

\titleformat{\subsection}
  {\normalfont\bfseries\normalsize}
  {\thesubsection}{1em}{}
\titlespacing*{\subsection}{0pt}{10pt}{4pt}

\usepackage[none]{hyphenat}
\sloppy

